%% SCRAP: experiments/02-experiments/L8_WIRING_IMPLEMENTATION_SUMMARY
%% SOURCE: docs/working/experiments/02-experiments/L8_WIRING_IMPLEMENTATION_SUMMARY.md
%% STATUS: HISTORICAL
%% FITS: experiments/ch-l8
%% EDITORIAL: lifted — prose rewritten to press voice

\section{L8 Control-Wiring Implementation Summary}
\label{sec:l8-wiring-summary}

The L8 Jacquard mode selector was 90\% complete before this implementation:
the state machine computed a 4-bit mode and set Boolean flags in
\texttt{vm->ssm\_config}, but the execution paths of the four target loops
did not consult those flags. This work completed the final ``last mile'' by
adding conditional gates at each loop's execution point.

\subsection{Changes Made}

Four source locations received new gates. The implementation order followed
increasing complexity: L3 (single-function internal gate), L5 (medium
complexity, preserve previous value), L6 (symmetric to L5), L2 (call-site
gate in the hot execution path).

\paragraph{L3 — Linear Decay} Gate added to
\texttt{physics\_metadata\_apply\_linear\_decay()} in
\texttt{src/physics\_metadata.c}~line~176, placed after the frozen-word and
minimum-interval checks. When L3 is disabled, words retain heat without decay.

\paragraph{L5 — Window Width Inference} Gate added to Phase~2C of
\texttt{inference\_engine\_run()} in \texttt{src/inference\_engine.c}~line~538.
Previous \texttt{adaptive\_window\_width} is preserved when L5 is off (static
configuration).

\paragraph{L6 — Decay Slope Inference} Gate added to Phase~2D of
\texttt{inference\_engine\_run()} in \texttt{src/inference\_engine.c}~line~563.
Previous \texttt{adaptive\_decay\_slope} is preserved when L6 is off.

\paragraph{L2 — Rolling Window} Gate added at the call site in
\texttt{execute\_colon\_word()} in \texttt{src/vm.c}~line~1648. No execution
history is recorded when L2 is off.

\subsection{Mode Mapping}

The 16 valid L8 modes encode all combinations of the four controlled loops:

\begin{center}
\begin{tabular}{clccccl}
\toprule
Mode & L2 & L3 & L5 & L6 & Description \\
\midrule
C0  & -- & -- & -- & -- & Minimal overhead (startup default) \\
C4  & -- & \checkmark & -- & -- & Decay only (top-5\% candidate) \\
C7  & -- & \checkmark & \checkmark & \checkmark & Full inference without history \\
C9  & \checkmark & -- & -- & \checkmark & History + slope (top-5\% candidate) \\
C11 & \checkmark & -- & \checkmark & \checkmark & History + full inference \\
C12 & \checkmark & \checkmark & -- & -- & DoE-optimal baseline \\
C15 & \checkmark & \checkmark & \checkmark & \checkmark & All loops (maximum adaptivity) \\
\bottomrule
\end{tabular}
\end{center}

Modes C1--C3, C5--C6, C8, C10, C13--C14 represent intermediate combinations not
listed above.

\subsection{Test Results}

\begin{lstlisting}[language=bash]
make clean && make fastest
# Zero warnings, zero errors

make test
# Total: 782  Passed: 731  Failed: 0  Skipped: 49  Errors: 0
\end{lstlisting}

All 731 implemented tests pass with the default mode C0 (all loops off).

\subsection{Heartbeat Control Cycle}

Each heartbeat tick ($\sim$1~ms) executes in four steps:

\begin{enumerate}
  \item \textbf{Metrics collection} — entropy, CV, temporal decay, and stability
        score are computed from the rolling window and inference engine outputs.
  \item \textbf{Mode selection} — \texttt{ssm\_l8\_update()} evaluates the four
        metrics against thresholds and computes a 4-bit mode with 5-tick
        hysteresis to suppress rapid flapping.
  \item \textbf{Configuration application} — \texttt{ssm\_apply\_mode()} decodes
        the mode bits into the L2/L3/L5/L6 Boolean flags.
  \item \textbf{Runtime enforcement} — the execution-path gates read the flags
        at the next loop iteration.
\end{enumerate}

\subsection{Planned Extensions}

Three enhancements are scoped but not yet implemented:

\begin{itemize}
  \item \textbf{FORTH API} — \texttt{L8-MODE!}, \texttt{L8-MODE@}, and
        \texttt{L8-AUTO} words to support manual mode forcing for DoE validation
        and benchmarking.
  \item \textbf{Heartbeat CSV column} — add \texttt{l8\_mode} to the per-tick
        export so R analysis can correlate mode transitions with performance
        metrics.
  \item \textbf{DoE re-validation} — run the full 128-configuration experiment
        with L8 wiring active and confirm 0\% algorithmic variance is maintained
        under dynamic mode switching.
\end{itemize}
